% Part: turing-machines
% Chapter: machines-computations
% Section: introduction

\documentclass[../../../include/open-logic-section]{subfiles}

\begin{document}

% SEGMENT OLP-0254-S01

\olfileid{tur}{mac}{int}
\olsection{அறிமுகம்}

எடுத்துக்காட்டாக, $\Nat$ இலிருந்து $\Nat$ க்கான ஒரு சார்பு
\emph{கணிக்கத்தக்கது} என்பதன் பொருள் என்ன? ஒரு சார்பை ஒரு டியூரிங்
பொறியால் கணிக்க முடிந்தால் அது கணிக்கத்தக்கது என்பதே முதல்
விடைகளில் ஒன்றும் மிகவும் அறியப்பட்ட விடையும் ஆகும். இக்கருத்தை
ஆலன் டியூரிங் 1936 இல் வகுத்தார். டியூரிங் பொறிகள்
\emph{கணிப்பின் ஒரு மாதிரிக்கு} எடுத்துக்காட்டு---``கணிப்புச்
செயல்முறை'' என்ற கருத்தை வரையறுக்கும் கணிதத் துல்லியமான வழி அவை.
அது துல்லியமாக எதைக் குறிக்கிறது என்பது விவாதத்துக்குரியது;
ஆனால் கணிப்புச் செயல்முறைகளைக் குறிப்பிடுவதற்கான ஒரு வழி டியூரிங்
பொறிகள் என்பது பரவலாக ஏற்கப்படுகிறது. ``டியூரிங் பொறி'' என்ற சொல்,
அசையும் உறுப்புகள் கொண்ட ஓர் இயற்பொருள் பொறியின் படிமத்தை
நினைவூட்டினாலும், கண்டிப்பாகக் கூறினால் டியூரிங் பொறி முற்றிலும்
கணிதக் கட்டமைப்பு; எனவே அது கணிப்புச் செயல்முறை என்ற கருத்தின்
இலட்சிய வடிவம். எடுத்துக்காட்டாக, டியூரிங் பொறியின் நேரத் தேவைக்கும்
நினைவகத் தேவைக்கும் நாம் எந்தக் கட்டுப்பாடும் விதிப்பதில்லை:
அண்டத்திலுள்ள அணுக்களைவிட அதிகமான சேமிப்பிடமோ படிகளோ ஒரு
கணிப்புக்குத் தேவைப்பட்டாலும்கூட, டியூரிங் பொறிகள் அதைக் கணிக்கலாம்.

\begin{explain}
ஒரு சிறப்புவகை கற்பனைப் பொறிமுறைக்கான நிரலாக டியூரிங் பொறியைக்
கருதுவதே சிறந்ததாக இருக்கலாம். இப்பொறிமுறை ஒரு \emph{நாடாவையும்}
ஒரு \emph{வாசி-எழுது தலையையும்} கொண்டது. நமது டியூரிங் பொறிப்
பதிப்பில், நாடா ஒரு திசையில் (வலப்புறமாக) முடிவிலாதது; அது
\emph{கட்டங்களாகப்} பிரிக்கப்பட்டுள்ளது. ஒவ்வொரு கட்டத்திலும் ஒரு
முடிவுறு \emph{குறியெழுத்துத் தொகுதியிலிருந்து} ஒரு குறி இருக்கலாம்.
இத்தகைய குறியெழுத்துத் தொகுதிகள் எத்தனை வேறுபட்ட குறிகளையும்
கொண்டிருக்கலாம்; ஆனால் முக்கியமாக மூன்றே நமக்குப் போதும்:
$\TMendtape$, $\TMblank$, $\TMstroke$. பொறிமுறை தொடங்கப்படும்போது,
நாடாவின் இடக்கோடிக் கட்டத்தில் $\TMendtape$ உள்ளது; முடிவுறு
எண்ணிக்கையிலான சில கட்டங்களில் \emph{உள்ளீடு} உள்ளது; மற்ற
ஒவ்வொரு கட்டத்திலும் $\TMblank$ இருப்பதால் நாடா வெறுமையாக உள்ளது.
எந்த நேரத்திலும், பொறிமுறை முடிவுறு எண்ணிக்கையிலான
\emph{நிலைகளில்} ஒன்றில் உள்ளது. தொடக்கத்தில், தலை இடக்கோடிக்
கட்டத்தை ஒரு குறிப்பிடப்பட்ட \emph{தொடக்க நிலையில்} வருடுகிறது.
பொறிமுறையின் இயக்கத்தில் ஒவ்வொரு படியிலும், அப்போது வருடப்படும்
கட்டத்தின் உள்ளடக்கமும் பொறிமுறை இருக்கும் நிலையும் டியூரிங் பொறி
நிரலும் அடுத்து என்ன நடக்கும் என்பதைத் தீர்மானிக்கின்றன. டியூரிங்
பொறி நிரல், ஒரு நிலை~$q$ ஐயும் ஒரு குறி~$\sigma$ ஐயும் உள்ளீடாக
ஏற்று $\tuple{q', \sigma', D}$ என்ற மும்மையை வெளியிடும் ஒரு
பகுதிச் சார்பால் கொடுக்கப்படுகிறது. பொறிமுறை நிலை $q$ இல் இருந்து
குறி $\sigma$ ஐ வாசிக்கும் போதெல்லாம், நடப்புக் கட்டத்திலுள்ள
குறியை $\sigma'$ ஆல் மாற்றுகிறது; $D$ ஆனது முறையே
$\TMleft$, $\TMright$, $\TMstay$ என்பதற்கேற்ப தலை இடப்புறம்,
வலப்புறம் நகர்கிறது அல்லது அங்கேயே நிற்கிறது; பொறிமுறை
நிலை~$q'$ குச் செல்கிறது.

எடுத்துக்காட்டாக, \olref{fig:tm} இலுள்ள சூழலைக் கருதுக.
\begin{figure}
  \olasset{\olpath/assets/diagrams/turing-machine.tikz}
  \caption{தன் நிரலை இயக்கும் ஒரு டியூரிங் பொறி.}
  \ollabel{fig:tm}
\end{figure}
டியூரிங் பொறியின் நாடாவில் தெரியும் பகுதி, இடக்கோடிக் கட்டத்தில்
நாடா-முடிவுக் குறி $\TMendtape$ ஐக் கொண்டுள்ளது; அதன்பின் மூன்று
$1$-கள், ஒரு $0$, மேலும் நான்கு $1$-கள் உள்ளன. தலை இடப்புறத்திலிருந்து
மூன்றாவது கட்டத்தை வாசிக்கிறது; அக்கட்டத்தில் ஒரு~$1$ உள்ளது;
தலை நிலை~$q_1$ இல் உள்ளது---``பொறி நிலை~$q_1$ இல் ஒரு
$1$ ஐ வாசிக்கிறது'' என்கிறோம். உள்ளீடு $\tuple{q_1, 1}$ க்கு
டியூரிங் பொறி நிரல் மும்மை $\tuple{q_2, 0, \TMstay}$ ஐத்
திருப்பித்தந்தால், பொறி இப்போது மூன்றாவது கட்டத்திலுள்ள~$1$ ஐ
ஒரு~$0$ ஆல் மாற்றி, வாசி-எழுது தலையை இருந்த இடத்திலேயே விட்டுவிட்டு,
நிலை~$q_2$ குச் செல்லும். பிறகு உள்ளீடு $\tuple{q_2, 0}$ க்கு
நிரல் $\tuple{q_3, 0, \TMright}$ ஐத் திருப்பித்தந்தால், பொறி
இப்போது~$0$ ஐ மற்றொரு~$0$ ஆல் மேலெழுதும் (இதனால் வாசி-எழுது
தலையின்கீழுள்ள நாடாவின் உள்ளடக்கம் மாறாமல் இருக்கும்), ஒரு கட்டம்
வலப்புறம் நகர்ந்து நிலை~$q_3$ ஐ அடையும். இவ்வாறே தொடரும்.

நிலை $q_n$, குறி $\sigma$ ஆகியவற்றை எதிர்கொள்ளும்போது
$\tuple{q_n, \sigma}$ கான எந்தக் கட்டளையும் இல்லாவிட்டால்,
அதாவது உள்ளீடு $\tuple{q_n,\sigma}$ கான நிலைமாற்றுச் சார்பு
வரையறுக்கப்படவில்லை எனில், பொறி \emph{நிற்கிறது} என்கிறோம்.
வேறுவிதமாகக் கூறினால், பொறி நிறைவேற்ற வேண்டிய கட்டளை எதுவும்
இல்லை; அப்புள்ளியில் அது செயல்படுவதை நிறுத்துகிறது. நிற்றல் சில
நேரங்களில் ஒரு குறிப்பிட்ட நிற்றல் நிலை~$h$ ஆல் குறிக்கப்படுகிறது.
இது பின்னர் மேலும் விரிவாகக் காட்டப்படும்.
\end{explain}

\begin{digress}
``கணிக்கத்தக்க எண்கள் குறித்து'' என்ற டியூரிங்கின் கட்டுரையின்
சிறப்பு, அவர் ஒரு முறைசார்ந்த வரையறையை மட்டுமல்லாமல், அவ்வரையறை
கணிப்புத்தன்மையின் உள்ளுணர்வுக் கருத்தைக் கைப்பற்றுகிறது என்பதற்கான
வாதத்தையும் முன்வைப்பதுதான். வரையறையிலிருந்து, ஒரு டியூரிங்
பொறியால் கணிக்கத்தக்க எந்தச் சார்பும் உள்ளுணர்வு நோக்கில்
கணிக்கத்தக்கது என்பது தெளிவாக இருக்க வேண்டும். மறுதலை உண்மை,
அதாவது இயல்பாகக் கணிக்கத்தக்கது என்று நாம் கருதும் எந்தச் சார்பும்
இத்தகைய ஒரு பொறியால் கணிக்கத்தக்கது என்பதற்கு டியூரிங் மூன்றுவகை
வாதங்களை வழங்குகிறார். அவை (டியூரிங்கின் சொற்களில்):
\begin{enumerate}
\item உள்ளுணர்வுக்கு நேரடி முறையீடு.
\item இரு வரையறைகளின் சமானத்தன்மைக்கான நிரூபணம்
  (புதிய வரையறை மேலும் வலுவான உள்ளுணர்வுக் கவர்ச்சியைக் கொண்டிருந்தால்).
\item கணிக்கத்தக்க எண்களின் பெரும் வகுப்புகளுக்கு எடுத்துக்காட்டுகளை வழங்குதல்.
\end{enumerate}
``கொள்கையளவில்'' கணிப்புத்தன்மை என்ற கருத்தை, அதாவது நேரம்,
இடம் ஆகியவற்றின் நடைமுறைக் கட்டுப்பாடுகளைக் கருத்தில் கொள்ளாமல்
வரையறுக்க முயல்வதே நமது குறிக்கோள். கணிப்புத்தன்மையின் மிகப்
பரந்த வரையறையை அமைத்தபின், வரம்பிட்ட வளங்களைக் கொண்ட கணிப்பையும்
கருதலாம் என்பது உண்மையே; ``கணிப்புச் சிக்கற்பாடு'' எனப்படும்
பாடப்புலத்தின் மையம் இதுவே.
\end{digress}

\begin{history}
ஆலன் டியூரிங் 1936 இல் டியூரிங் பொறிகளைக் கண்டுபிடித்தார்.
அப்போது அவரது ஆர்வம் முதல்தர தருக்கத்தின் தீர்மானிக்கத்தன்மையாக
இருந்தபோதிலும், அக்கட்டுரை கணினி வடிவமைப்பின் அடித்தளங்கள் குறித்த
ஒரு முடிவான கட்டுரை என விவரிக்கப்பட்டுள்ளது. அக்கட்டுரையில்,
கணிக்கத்தக்க மெய்யெண்கள், அதாவது கணிக்கத்தக்க தசம விரிவாக்கங்களைக்
கொண்ட மெய்யெண்கள் மீது டியூரிங் கவனம் செலுத்துகிறார்; ஆனால் அவரது
கருத்துகளை இயல் எண்கள்மீதான கணிக்கத்தக்க சார்புகளுக்கும் பிறவற்றுக்கும்
ஏற்ப மாற்றுவது கடினமல்ல என்று குறிப்பிடுகிறார். முதல் செயல்படும்
பொதுநோக்குக் கணினி 1941 இல் (ஜெர்மானியரான கொன்ராட் சூசே தம்
பெற்றோரின் வரவேற்பறையில் கட்டியது) உருவாக்கப்படுவதற்கு முழுமையாக
ஐந்து ஆண்டுகளுக்கு முன்பே இது நிகழ்ந்தது. டியூரிங்கும் பிளெட்ச்லி
பார்க்கிலிருந்த அவரது உடன் பணியாளர்களும் மறைக்குறி முறிக்கும்
கொலோசஸை (1943) கட்டுவதற்கு ஏழு ஆண்டுகளுக்கு முன்பு; அமெரிக்க
ENIAC (1945) க்கு ஒன்பது ஆண்டுகளுக்கு முன்பு; முதல் பிரித்தானியப்
பொதுநோக்குக் கணினியான மான்செஸ்டர் சிற்றளவு சோதனைப் பொறி
மான்செஸ்டரில் (1948) கட்டப்படுவதற்குப் பன்னிரண்டு ஆண்டுகளுக்கு
முன்பு; அமெரிக்கர்கள் BINAC ஐ முதலில் சோதித்ததற்கு (1949)
பதின்மூன்று ஆண்டுகளுக்கு முன்பு. முதல் சேமித்த-நிரல் கணினி என்ற
சிறப்பு மான்செஸ்டர் SSEM க்கு உண்டு---முந்தைய பொறிகள் ஒவ்வொரு
புதிய பணிக்கும் கையால் மீண்டும் கம்பியிடப்பட வேண்டியிருந்தன.
\end{history}

\end{document}
